\documentclass[12pt,a4paper]{article}
\usepackage[utf8]{inputenc}
\usepackage{babel}
\usepackage[T1]{fontenc}
\usepackage{geometry}
\geometry{margin=2.5cm}
\usepackage{graphicx}
\usepackage{booktabs}
\usepackage{amsmath}
\usepackage{hyperref}
\usepackage{float}
\usepackage{caption}
\usepackage{subcaption}
\usepackage{xcolor}
\usepackage{parskip}
\usepackage{lmodern}

\hypersetup{
    colorlinks=true,
    linkcolor=blue,
    urlcolor=blue
}

\title{\textbf{Práctica Final:\\Análisis y Predicción de Series Temporales}}
\author{Gonzalo Cruz Gómez}
\date{Mayo 2026}

\begin{document}
\maketitle
\setcounter{tocdepth}{1}
\tableofcontents
\newpage

\section{Introducción}

El objetivo de esta práctica es predecir cuatro series temporales de naturaleza muy diferente para la competición Kaggle de la asignatura. Para cada una se ha seguido el mismo proceso: análisis exploratorio, test de estacionariedad, selección del modelo por RMSE en validación y predicción final reajustando sobre todos los datos disponibles.

El enfoque utilizado en todas las series es SARIMA/SARIMAX (Seasonal AutoRegressive Integrated Moving Average with eXogenous regressors). Para la serie B, que presenta doble estacionalidad (semanal y anual), se combina SARIMA con términos de Fourier anuales como regresores externos. En la serie C se añade también un regresor de tendencia lineal.

Las herramientas usadas son Python 3.10 con \texttt{pandas}, \texttt{statsmodels} y \texttt{matplotlib}.

\section{Descripción de los Datos}

\begin{table}[H]
\centering
\caption{Resumen de las series temporales}
\begin{tabular}{llcccc}
\toprule
\textbf{Serie} & \textbf{Variable} & \textbf{Frecuencia} & \textbf{Obs. train} & \textbf{Período train} & \textbf{Preds.} \\
\midrule
A & Precio bursátil     & Diaria (lab.) & 2025 & 2010--2018 & 10 \\
B & Nacimientos diarios & Diaria        & 3287 & 1994--2002 & 365 \\
C & Temp. global (°C)   & Mensual       & 1740 & 1880--2024 & 12 \\
D & Serie mensual       & Mensual       & 132  & 1980--1990 & 10 \\
\bottomrule
\end{tabular}
\end{table}

\section{Serie A - Precios Bursátiles Diarios}

\subsection{Análisis exploratorio}

La serie A contiene 2025 observaciones de precios bursátiles en días laborables (julio 2010 -- septiembre 2018). La tendencia alcista es clara: el precio pasa de unos 120 a más de 220.

\begin{figure}[H]
    \centering
    \includegraphics[width=\linewidth]{figures/A_serie.png}
    \caption{Serie A: precios bursátiles diarios (2010--2018)}
\end{figure}

\begin{figure}[H]
    \centering
    \includegraphics[width=\linewidth]{figures/A_descomposicion.png}
    \caption{Descomposición aditiva de la Serie A (período = 5 días laborables)}
\end{figure}

\subsection{Estacionariedad}

El test ADF en niveles da $p=0.74$, así que la serie no es estacionaria. Con la primera diferencia cae a $p<0.0001$, lo que confirma integración de orden 1. El ACF y PACF de la serie diferenciada no muestran autocorrelaciones relevantes, lo que es coherente con un precio que sigue aproximadamente un paseo aleatorio.

\begin{figure}[H]
    \centering
    \includegraphics[width=\linewidth]{figures/A_acf_pacf.png}
    \caption{ACF y PACF de la Serie A (primera diferencia)}
\end{figure}

\subsection{Selección del modelo}

Se realizó un grid search SARIMA con $m=5$ días laborables (estacionalidad semanal). El mejor modelo en validación (últimas 20 observaciones) fue \textbf{SARIMA$(2,0,1)(1,1,1)_5$} con RMSE$=6.16$.

Con $d=0$ y $D=1$, la diferenciación estacional estabiliza la serie sin perder la información de nivel. El término AR$(2)$ ordinario captura la dinámica de corto plazo (la correlación entre días consecutivos).

\begin{table}[H]
\centering
\caption{Selección SARIMA - Serie A (validación: últimas 20 obs.)}
\begin{tabular}{lcc}
\toprule
\textbf{Modelo} & \textbf{RMSE (val.)} & \textbf{AIC} \\
\midrule
\textbf{SARIMA$(2,0,1)(1,1,1)_5$} & \textbf{6.16} & 10378.55 \\
SARIMA$(2,0,1)(0,1,1)_5$           & 6.24 & 10349.53 \\
SARIMA$(0,0,2)(1,0,0)_5$           & 6.26 & 11746.16 \\
\bottomrule
\end{tabular}
\end{table}

\subsection{Análisis de residuos}

\begin{figure}[H]
    \centering
    \includegraphics[width=\linewidth]{figures/A_residuos.png}
    \caption{Análisis de residuos -- SARIMA$(2,0,1)(1,1,1)_5$ -- Serie A}
\end{figure}

\subsection{Predicciones}

\begin{figure}[H]
    \centering
    \includegraphics[width=\linewidth]{figures/A_predicciones.png}
    \caption{Serie A: predicciones SARIMA$(2,0,1)(1,1,1)_5$ para los 10 días de test}
\end{figure}

\section{Serie B - Nacimientos Diarios}

\subsection{Análisis exploratorio}

La serie B contiene 3287 observaciones diarias de nacimientos (1994-2002). Tiene dos estacionalidades claramente visibles: una semanal muy marcada (pocos nacimientos en fin de semana, las cesáreas programadas se concentran entre semana) y una anual (pico en agosto-septiembre, mínimo en enero, diferencia de unos 1000 nacimientos/día).

\begin{figure}[H]
    \centering
    \includegraphics[width=\linewidth]{figures/B_serie.png}
    \caption{Serie B: nacimientos diarios (1994--2002)}
\end{figure}

\begin{figure}[H]
    \centering
    \includegraphics[width=\linewidth]{figures/B_descomposicion.png}
    \caption{Descomposición de la Serie B (período = 7 días)}
\end{figure}

\subsection{Modelo: SARIMA con términos de Fourier}

La serie tiene dos estacionalidades distintas. SARIMA con $m=7$ captura directamente la componente semanal mediante diferenciación estacional y los términos AR/MA estacionales. Sin embargo, la componente anual no puede modelarse con un segundo periodo estacional en SARIMA estándar.

Para resolver esto se añaden como regresores externos \textbf{términos de Fourier anuales}: pares seno/coseno de la forma
\[
\sin\!\left(\frac{2\pi k\, t}{365.25}\right), \quad \cos\!\left(\frac{2\pi k\, t}{365.25}\right), \quad k = 1, \ldots, K
\]
donde $t$ es el día del año (1--365). Con $K=5$ se tienen 10 regresores que representan los primeros cinco armónicos de la variación anual. La ventaja de este enfoque es que estos regresores se calculan únicamente a partir del índice de fecha, sin usar los valores de la serie, así que no hay fuga de información al construirlos para 2003. El modelo SARIMAX con $m=7$ y Fourier $K=5$ separa entonces las dos fuentes de variación: el componente estocástico semanal (manejado por el SARIMA) y el patrón anual determinista (manejado por el Fourier).

\subsubsection*{Selección del modelo}

Se realizó un grid search con $m=7$ y $K=5$, entrenando con 1994-2001 y validando en 2002. El mejor modelo fue \textbf{SARIMA$(2,0,1)(1,1,1)_7$} con Fourier $K=5$ (RMSE$=710$ sobre el año 2002, entrenado con todos los años disponibles).

\begin{table}[H]
\centering
\caption{Top modelos - Serie B (validación: año 2002)}
\begin{tabular}{lccc}
\toprule
\textbf{Modelo} & \textbf{K} & \textbf{Ventana} & \textbf{RMSE (val.)} \\
\midrule
\textbf{SARIMA$(2,0,1)(1,1,1)_7$} & \textbf{5} & \textbf{all} & \textbf{710.10} \\
SARIMA$(2,0,0)(1,1,1)_7$          & 5 & all & 712.25 \\
SARIMA$(1,0,2)(1,1,1)_7$          & 5 & all & 712.63 \\
\bottomrule
\end{tabular}
\end{table}

\subsection{Análisis de residuos}

\begin{figure}[H]
    \centering
    \includegraphics[width=\linewidth]{figures/B_residuos.png}
    \caption{Análisis de residuos - SARIMA$(2,0,1)(1,1,1)_7$ + Fourier $K=5$ -- Serie B}
\end{figure}

\subsection{Predicciones}

\begin{figure}[H]
    \centering
    \includegraphics[width=\linewidth]{figures/B_predicciones.png}
    \caption{Serie B: predicciones para el año 2003 (SARIMA + Fourier $K=5$)}
\end{figure}

\section{Serie C - Temperatura Global Mensual}

\subsection{Análisis exploratorio}

La serie C tiene temperaturas medias globales mensuales desde 1880 hasta 2024 (1740 observaciones). Se observa claramente el calentamiento global como tendencia ascendente, especialmente acelerado desde los años 80, y una pequeña estacionalidad anual.

\begin{figure}[H]
    \centering
    \includegraphics[width=\linewidth]{figures/C_serie.png}
    \caption{Serie C: temperatura global mensual (°C), 1880--2024}
\end{figure}

\begin{figure}[H]
    \centering
    \includegraphics[width=\linewidth]{figures/C_descomposicion.png}
    \caption{Descomposición de la Serie C (período = 12 meses)}
\end{figure}

\subsection{Estacionariedad}

El ADF en niveles da $p=0.97$ (no estacionaria). Con diferenciación estacional de orden 12 baja a $p<0.0001$.

\subsection{Selección del modelo}

Se entrenó sobre los \textbf{últimos 40 años} (1985--2024, 480 observaciones) en vez de usar la serie completa desde 1880. La razón es que la tendencia de calentamiento reciente es mucho más pronunciada que la de finales del siglo XIX y principios del XX: los datos antiguos tienen una pendiente muy diferente a la actual y distorsionan las predicciones a corto plazo.

El mejor modelo fue \textbf{SARIMA$(1,0,0)(0,1,1)_{12}$} con un \textbf{regresor de tendencia lineal} (RMSE$=0.287$ en validación, últimos 24 meses). La tendencia lineal exógena modela el calentamiento global continuo de forma explícita; la diferenciación estacional $D=1$ elimina el patrón mensual.

\begin{table}[H]
\centering
\caption{Top modelos - Serie C (validación: 2023--2024)}
\begin{tabular}{lcc}
\toprule
\textbf{Modelo} & \textbf{Ventana} & \textbf{RMSE (val.)} \\
\midrule
\textbf{SARIMA$(1,0,0)(0,1,1)_{12}$ + tendencia} & \textbf{40 años} & \textbf{0.2870} \\
SARIMA$(0,0,1)(0,1,1)_{12}$ + tendencia          & 40 años & 0.2894 \\
SARIMA$(0,0,1)(0,1,1)_{12}$                       & 40 años & 0.2895 \\
\bottomrule
\end{tabular}
\end{table}

\subsection{Análisis de residuos}

\begin{figure}[H]
    \centering
    \includegraphics[width=\linewidth]{figures/C_residuos.png}
    \caption{Análisis de residuos - SARIMA$(1,0,0)(0,1,1)_{12}$ + tendencia -- Serie C}
\end{figure}

\subsection{Predicciones}

\begin{figure}[H]
    \centering
    \includegraphics[width=\linewidth]{figures/C_predicciones.png}
    \caption{Serie C: predicciones SARIMA$(1,0,0)(0,1,1)_{12}$ para 2025}
\end{figure}

\section{Serie D - Serie Mensual}

\subsection{Análisis exploratorio}

La serie D tiene 132 observaciones mensuales (1980--1990). Los valores crecen de forma sostenida (de $\approx 464$ a $\approx 2286$) con una estacionalidad anual clara. Probablemente son datos de consumo (gas o electricidad).

\begin{figure}[H]
    \centering
    \includegraphics[width=\linewidth]{figures/D_serie.png}
    \caption{Serie D: serie mensual (1980--1990)}
\end{figure}

\begin{figure}[H]
    \centering
    \includegraphics[width=\linewidth]{figures/D_descomposicion.png}
    \caption{Descomposición de la Serie D (período = 12 meses)}
\end{figure}

\subsection{Estacionariedad}

ADF en niveles: $p=0.92$ (no estacionaria). Con diferenciación estacional: $p=0.0003$.

\subsection{Selección del modelo}

El mejor modelo en validación (año 1990) fue \textbf{SARIMA$(0,0,0)(1,0,0)_{12}$} con RMSE$=186.24$. Este modelo solo tiene el componente AR estacional de orden 1: cada observación es una fracción del mismo mes del año anterior. Su simplicidad es su ventaja, ya que no sobre-ajusta en una serie de solo 132 observaciones.

\begin{table}[H]
\centering
\caption{Top modelos - Serie D (validación: año 1990)}
\begin{tabular}{lcc}
\toprule
\textbf{Modelo} & \textbf{RMSE (val.)} & \textbf{AIC} \\
\midrule
\textbf{SARIMA$(0,0,0)(1,0,0)_{12}$} & \textbf{186.24} & 1483.61 \\
SARIMA$(0,0,1)(1,0,0)_{12}$          & 193.89 & 1473.83 \\
SARIMA$(0,0,0)(0,1,1)_{12}$          & 193.98 & 1283.23 \\
\bottomrule
\end{tabular}
\end{table}

\subsection{Análisis de residuos}

\begin{figure}[H]
    \centering
    \includegraphics[width=\linewidth]{figures/D_residuos.png}
    \caption{Análisis de residuos - SARIMA$(0,0,0)(1,0,0)_{12}$ - Serie D}
\end{figure}

\subsection{Predicciones}

\begin{figure}[H]
    \centering
    \includegraphics[width=\linewidth]{figures/D_predicciones.png}
    \caption{Serie D: predicciones para enero--octubre de 1991}
\end{figure}

\newpage
\section{Investigación Adicional}
\label{sec:inv}

Además del análisis principal, exploré algunas cuestiones adicionales que aparecieron durante el trabajo.

\subsection{Serie A -- Heterocedasticidad y log-retornos}

La serie de precios es heterocedástica: cuando el precio es alto, también lo son las oscilaciones diarias. Aplicando el test de Levene dividiendo la serie en 4 bloques, el p-valor es prácticamente cero.

En finanzas lo habitual es trabajar con \textbf{log-retornos} ($r_t = \log P_t - \log P_{t-1}$), que son estacionarios y mucho más homocedásticos. Ajusté un ARIMA sobre los retornos y luego revertí la transformación para obtener precios. El modelo resultante fue ARIMA$(0,0,0)$ (ruido blanco), coherente con la hipótesis de eficiencia del mercado.

\begin{figure}[H]
    \centering
    \includegraphics[width=\linewidth]{figures/inv_A_heterocedasticidad.png}
    \caption{Serie A: desviación típica móvil de precios (izquierda) y log-precios (derecha). La transformación reduce bastante la heterocedasticidad.}
\end{figure}

\begin{figure}[H]
    \centering
    \includegraphics[width=\linewidth]{figures/inv_A_logretornos.png}
    \caption{Arriba: log-retornos (rachas de alta y baja volatilidad). Abajo: predicción en precios obtenida revirtiendo la transformación logarítmica.}
\end{figure}


\subsection{Serie D - Transformación Box-Cox}

La transformación Box-Cox elige automáticamente el exponente que mejor estabiliza la varianza. La serie D crece bastante así que pensé que podría haber heterocedasticidad, pero el test de Levene dio $p=0.31$: la varianza ya era constante. El lambda óptimo resultó $\lambda=0.35$, entre logaritmo y raíz cuadrada.

\begin{figure}[H]
    \centering
    \includegraphics[width=\linewidth]{figures/inv_D_boxcox.png}
    \caption{Desviación típica móvil antes y después de Box-Cox ($\lambda=0.35$). La mejora es pequeña porque la serie ya era bastante homocedástica.}
\end{figure}

Ajusté un SARIMA sobre la serie transformada y el RMSE en validación fue de 264, peor que el modelo directo (186). La transformación no ayuda porque el problema no era heterocedasticidad sino tendencia, que ya maneja el SARIMA con sus componentes AR.

\section{Resultados y Conclusiones}

\begin{table}[H]
\centering
\caption{Modelos finales y métricas de validación}
\begin{tabular}{llcc}
\toprule
\textbf{Serie} & \textbf{Modelo} & \textbf{RMSE (val.)} & \textbf{Obs. train} \\
\midrule
A & SARIMA$(2,0,1)(1,1,1)_5$                        & 6.16  & 2025 \\
B & SARIMA$(2,0,1)(1,1,1)_7$ + Fourier $K=5$         & 710.10 & 3287 \\
C & SARIMA$(1,0,0)(0,1,1)_{12}$ + tendencia (40 años) & 0.287 & 480 \\
D & SARIMA$(0,0,0)(1,0,0)_{12}$                       & 186.24 & 132 \\
\bottomrule
\end{tabular}
\end{table}

Las principales conclusiones de la práctica son:

\begin{itemize}
    \item \textbf{Serie A}: La diferenciación estacional $D=1$ con $m=5$ captura el patrón de semana laboral. Los parámetros AR$(2)$, MA$(1)$ y los estacionales AR$(1)$, MA$(1)$ añaden flexibilidad para modelar la autocorrelación de corto y largo plazo.

    \item \textbf{Serie B}: La doble estacionalidad se puede manejar combinando SARIMA (para la componente semanal) con términos de Fourier como regresores externos (para la componente anual). El uso de Fourier es especialmente útil cuando la serie tiene dos frecuencias estacionales que no pueden capturarse con un solo periodo $m$.

    \item \textbf{Serie C}: Limitar los datos de entrenamiento a los últimos 40 años es una decisión importante: los datos de 1880--1984 tienen una pendiente de calentamiento muy diferente a la actual y añadir un regresor de tendencia lineal junto con la ventana temporal adecuada mejora sustancialmente las predicciones.

    \item \textbf{Serie D}: Con solo 132 observaciones, los modelos simples generalizan mejor. SARIMA$(0,0,0)(1,0,0)_{12}$ con un único parámetro estacional fue el más robusto en validación.
\end{itemize}

\subsection*{Trabajo futuro}

Para la serie B, una línea interesante sería usar el \textbf{HistGradientBoosting} como modelo principal. Los resultados en validación sugieren que los lags anuales combinados con features de calendario pueden capturar mejor las interacciones no lineales entre el día de la semana, el mes y el nivel del año anterior. Con una validación rolling sobre varios años y una búsqueda de hiperparámetros más exhaustiva, podría mejorar sustancialmente el RMSE.

\end{document}
